\section*{Problems}

\textbf{3.1.} For integer $n = 1, 2, 3, \dots$, let $\text{rem}(n, 3)$ denote the remainder obtained after dividing $n$ by 3. For example, $\text{rem}(5, 3) = 2$ and $\text{rem}(6, 3) = 0$. For $n \ge 1$, let $A_n$ be the closed interval
\[
[\text{rem}(n, 3), \ 5 + e^{-n}] :
\]
\begin{enumerate}
    \item Find $\limsup_n A_n$ and $\liminf_n A_n$.
    \item Let $F(x)$ be a Stieltjes measure function defined as follows:
    \[
    F(x) = \begin{cases} 
    0 & \text{if } x < 0 \\ 
    x/4 & \text{if } 0 \le x < 2 \\ 
    1 & \text{if } 2 \le x. 
    \end{cases}
    \]
    We use $F(x)$ to generate a Lebesgue--Stieltjes measure $P$. Compute the probabilities $P(\limsup_n A_n)$ and $P(\liminf_n A_n)$.
\end{enumerate}

\textbf{3.2.} Let $P$ be a Lebesgue--Stieltjes measure on $\mathbb{R}$ corresponding to a continuous distribution function, i.e., $P$ is a set function on $\mathcal{B}(\mathbb{R})$ and there exists a continuous function $F(x)$ such that $P([a, b]) = F(b) - F(a)$ for all $a < b$, and $P(\mathbb{R}) = 1$:
\begin{enumerate}
    \item Show that $P(A) = 0$ for all countable subsets $A$ of $\mathbb{R}$.
    \item If $P(A) = 1$, can we deduce that $A$ is a dense subset of $\mathbb{R}$?
\end{enumerate}

\textbf{3.3.} For $i = 1, 2, 3, \dots$, let $p_i$ be a nonnegative real number such that $\sum_{i=1}^\infty p_i$ is finite. We define a function $\mu$ on $2^{\mathbb{N}}$ by

\[
\mu(A) \triangleq \begin{cases} 
\sum_{i \in A} p_i & \text{if } A \text{ has finite cardinality} \\ 
\infty & \text{if } A \text{ is an infinite set,} 
\end{cases}
\]
for $A \subseteq \mathbb{N}$:
\begin{enumerate}[label=(\alph*)]
    \item Determine whether $\mu$ is finitely additive.
    \item Determine whether $\mu$ is $\sigma$-additive.
\end{enumerate}

\textbf{3.4.} Suppose $F_1(x)$ and $F_2(x)$ are two Stieltjes measure functions defined on $\mathbb{R}$, and let $\mu_1$ and $\mu_2$ be the corresponding Lebesgue--Stieltjes measures. Show that if $\mu_1$ and $\mu_2$ are the same set functions, then $F_1(x)$ and $F_2(x)$ differ by an additive constant.

\textbf{3.5.} Let $F(x)$ be a Stieltjes measure function that corresponds to a probability measure. Show that the set of points at which $F(x)$ has a jump discontinuity is at most countably infinite.

\textbf{3.6.} Consider a sample space $\Omega = \{a, b, c, d, e, f\}$ and $\sigma$-algebra $\mathscr{F} = 2^\Omega$. We define two different probability measures $P$ and $Q$ on this measure space by $P(\omega) = 1/6$ for all $\omega \in \Omega$, and
\[
Q(b) = Q(e) = 1/2, \quad Q(a) = Q(c) = Q(d) = Q(f) = 0.
\]
Let $\mathscr{C}$ denote the collection of sets
\[
\{\{a, b, c\}, \{b, c, d\}, \{c, d, e\}, \{d, e, f\}, \{e, f, a\}, \{f, a, b\}\} :
\]
\begin{enumerate}[label=(\alph*)]
    \item Prove that the $\sigma$-algebra generated by $\mathscr{C}$ is the same as $\mathscr{F}$.
    \item Show that $P(A) = Q(A)$ for all $A \in \mathscr{C}$.
    \item Does it violate the uniqueness of measure extension in Theorem 3.8?
\end{enumerate}

\textbf{3.7.} (An alternate definition of $\lambda$-system) Show that the following three conditions are equivalent to the definition of $\lambda$-system in Definition 3.10:
\begin{enumerate}
    \item $\Omega \in \mathcal{L}$.
    \item If $A$ and $B$ are sets in $\mathcal{L}$ and $A \subseteq B$, then $B \setminus A \in \mathcal{L}$.
    \item If $A_1, A_2, A_3, \dots$ is a sequence of \textit{increasing} sets in $\mathcal{L}$, then $\cup_{i=1}^{\infty} A_i \in \mathcal{L}$.
\end{enumerate}

\textbf{3.8.} Show that a $\lambda$-system is a $\sigma$-algebra if and only if it is a $\pi$-system.

\textbf{3.9.} (Proof of the $\pi$-$\lambda$ theorem) Let $\mathcal{P}$ be a non-empty $\pi$-system contained in a $\lambda$-system $\mathcal{L}$:



\begin{enumerate}[label=(\alph*)]
    \item Show that the intersection of all $\lambda$-systems included in $\mathcal{L}$ that contain $\mathcal{P}$ is a $\lambda$-system. Denote this $\lambda$-system by $\mathcal{L}_0$.
    \item Prove that $\mathcal{L}_1 \triangleq \{A \in \mathcal{L}_0 : A \cap B \in \mathcal{L}_0 \text{ for all } B \in \mathcal{P}\}$ is equal to $\mathcal{L}_0$ by showing $\mathcal{L}_1 is a $\lambda$-system.
    \item Prove that $\mathcal{L}_2 \triangleq \{A \in \mathcal{L}_1 : A \cap C \in \mathcal{L}_0 \text{ for all } C \in \mathcal{L}_0\}$ is equal to $\mathcal{L}_0$, by showing $\mathcal{L}_2$ is a $\lambda$-system.
    \item Use the previous exercise to show that $\mathcal{L}_0$ is a $\sigma$-algebra containing $\mathcal{P}$.
\end{enumerate}